\chapter{Tâche 3 — Rétro-ingénierie pour la compilation avec un SDK}
\label{chap:t3}
% Entreprise : Beemo Technologie

\section{Contexte et dette technique initiale}
% TODO : un programme doit être compilé avec un \gls{sdk} donné, mais la
% compilation échoue / n'est pas documentée. C'est un cas typique de dette
% (build cassé, savoir perdu, dépendance opaque). Enjeu business : sans
% build, pas de livraison.
\emph{[À rédiger.]}

\section{Objectif de la mission}
% TODO : rendre la compilation fonctionnelle avec le \gls{sdk}, ce qui a
% nécessité une démarche de \gls{retro-ingenierie}.
\emph{[À rédiger.]}

\section{Démarche de rétro-ingénierie}
% TODO : c'est le cœur de cette tâche — détailler la méthode :
% - analyse des artefacts (binaires, symboles, \gls{abi}/\gls{api}, headers)
% - reconstitution de la \gls{toolchain} / des dépendances manquantes
% - outils utilisés (objdump, nm, strings, ldd, débogueur... à adapter)
% - itérations d'essais-erreurs
\emph{[À rédiger.]}

\section{Réalisation}
% TODO : extrait de configuration de build / commande / patch ayant
% débloqué la compilation (commenté).
% \begin{lstlisting}[language=bash,caption={Étape clé de la compilation reconstituée.}]
% ...
% \end{lstlisting}
\emph{[À rédiger.]}

\section{Difficultés rencontrées}
% TODO : absence de doc, incompatibilités de versions, opacité du SDK.
\emph{[À rédiger.]}

\section{Résultats et bénéfice business}
% TODO : build reproductible et documenté = dette résorbée, déblocage du
% produit, capitalisation du savoir pour l'équipe.
\emph{[À rédiger.]}
